\documentclass[11pt,aspectratio=169]{beamer}
\usepackage{hp-dfr-slides}
\graphicspath{{../../../assets/}}

\title{From PINNs to a quantity of interest}
\subtitle{Goal-Oriented Deep Fourier Residual methods}
\author{Carlos Vera-Ciro}
\institute{Goal-Oriented DFR series \textbar\ Part 1}
\date{}

\begin{document}

\begin{frame}[plain,noframenumbering]
  \titlepage
\end{frame}

% =============================================================================
\begin{frame}{About this series}
  Four short decks, building the idea and then testing it honestly.
  \medskip
  \begin{enumerate}\itemsep0.55em
    \item[\important{1.}] \important{The problem} (this deck): why the residual
          norm matters, and why global error is often the wrong target.
    \item[\important{2.}] \important{The method}: goal-oriented DFR, a primal and
          an adjoint network trained together.
    \item[\important{3.}] \important{The theory}: the error representation and the
          bound, with proofs.
    \item[\important{4.}] \important{Results and limits}: what the method buys at
          equal compute, and where it does not help.
  \end{enumerate}
  \takeaway{Audience: anyone comfortable with PDEs and neural networks. No prior
  DFR needed.}
\end{frame}

% =============================================================================
\section{Solving PDEs with networks}
% =============================================================================

\begin{frame}{Solving a PDE with a neural network}
  We want the \important{exact solution} $u^*$: $-\Delta u^* = f$, with $u^* = 0$
  on $\partial\Omega$. We never have it, and every symbol below is named against it.
  \begin{itemize}\itemsep0.35em
    \item Represent the unknown by a network $u_\theta$ with parameters $\theta$.
    \item Choose $\theta$ so $u_\theta$ satisfies the equation as closely as
          possible, where ``closely'' is fixed by a \important{loss function}.
    \item The choice of loss decides how faithful the answer is.
  \end{itemize}
  \smallskip
  \begin{center}
  \begin{tikzpicture}[node distance=5mm and 9mm]
    \node[accentnode, minimum height=7mm] (net) {network $u_\theta$};
    \node[opnode, minimum height=7mm, right=of net] (res) {residual $f + \Delta u_\theta$};
    \node[opnode, minimum height=7mm, right=of res] (loss) {loss};
    \draw[localarrow] (net) -- (res);
    \draw[localarrow] (res) -- (loss);
    \coordinate (rtn) at ($(loss.south)+(0,-0.45)$);
    \draw[dashedarrow] (loss.south) -- (rtn) -| (net.south);
    \node[caption, below] at ($(rtn)!0.5!(rtn -| net.south)$) {update $\theta$};
  \end{tikzpicture}
  \end{center}
  \takeaway{Two objects to pin down: the \emph{network} $u_\theta$ (next two
  slides), then the \emph{loss} that scores it (the rest of this deck).}
\end{frame}

% -----------------------------------------------------------------------------
\begin{frame}{What the network is: a function of the coordinate}
  The \important{trainable network} is a neural field $\tilde{u}_\theta : \R^d \to \R$;
  every parameter $\theta$ lives in it. The \important{solution} we use is
  $u_\theta = \varphi\,\tilde{u}_\theta$, this network times a \emph{fixed} cutoff
  $\varphi$ (built on the next slide).
  \medskip
  \begin{columns}[c]
    \begin{column}{0.46\textwidth}
      \centering
      \begin{tikzpicture}[>=Latex]
        \foreach \y [count=\i] in {0.45,-0.45} \node[inneuron] (I\i) at (0,\y) {};
        \foreach \y [count=\i] in {1.15,0.45,-0.45,-1.15} \node[neuron] (A\i) at (1.5,\y) {};
        \foreach \y [count=\i] in {1.15,0.45,-0.45,-1.15} \node[neuron] (B\i) at (3.0,\y) {};
        \node[outneuron] (O) at (4.5,0) {};
        \foreach \i in {1,2} \foreach \j in {1,2,3,4} \draw[netedge] (I\i.east) -- (A\j.west);
        \foreach \i in {1,2,3,4} \foreach \j in {1,2,3,4} \draw[netedge] (A\i.east) -- (B\j.west);
        \foreach \i in {1,2,3,4} \draw[netedge] (B\i.east) -- (O.west);
        \node[left=2pt of I1] {\small $x_1$};
        \node[left=2pt of I2] {\small $x_2$};
        \node[right=2pt of O] {\small $\tilde{u}_\theta(x)$};
        \node[caption, align=center] at (2.25,-1.85) {the network $\tilde{u}_\theta$\\ 2D example ($d = 2$)};
      \end{tikzpicture}
    \end{column}
    \begin{column}{0.52\textwidth}
      \begin{itemize}\itemsep0.5em
        \item \important{Inputs:} the coordinate $x$, $d$ of them ($1$ in 1D, $2$
              in 2D). \emph{Not} one input per grid point.
        \item \important{Output:} one scalar, the raw value $\tilde{u}_\theta(x)$.
        \item The \emph{same} network is queried at every quadrature point ($150$
              in 1D, $32 \times 32$ in 2D); those are not extra inputs.
      \end{itemize}
    \end{column}
  \end{columns}
  \takeaway{``Degrees of freedom'' later means parameter count, not a mesh: e.g.
  $105$ for a $2 \to 8 \to 8 \to 1$ net.}
\end{frame}

% -----------------------------------------------------------------------------
\begin{frame}{Boundary conditions, built in by construction}
  The problem asks for $u^* = 0$ only \emph{on} the boundary $\partial\Omega$;
  inside, $u^*$ is the unknown we solve for. We enforce this exactly by
  construction, with no loss penalty to tune.
  \smallskip
  \begin{columns}[c]
    \begin{column}{0.60\textwidth}
      \begin{itemize}\itemsep0.4em
        \item $\tilde{u}_\theta$ is the trainable \important{network} from the last
              slide; on its own it is nonzero on $\partial\Omega$.
        \item $\varphi$ is a \important{fixed, known} function of the \emph{same}
              coordinate $x$ that feeds the network (never trained): zero on
              $\partial\Omega$, positive inside. On $[0,1]$, $\varphi(x) = x(1-x)$;
              on the square, $x_1(1-x_1)\,x_2(1-x_2)$.
        \item The product $u_\theta = \varphi\,\tilde{u}_\theta$ is then zero on
              $\partial\Omega$ (as $\varphi = 0$) and free inside (as $\varphi > 0$).
      \end{itemize}
    \end{column}
    \begin{column}{0.38\textwidth}
      \centering
      \begin{tikzpicture}[>=Latex]
        \draw[->] (-0.15,0) -- (3.5,0) node[right, caption] {$x$};
        \draw[->] (0,-0.15) -- (0,2.15) node[left, caption] {$\varphi$};
        \draw[thick, accentdark, domain=0:3, samples=80, smooth]
          plot (\x, {6.4*(\x/3)*(1-\x/3)});
        \fill[danger] (0,0) circle (1.6pt);
        \fill[danger] (3,0) circle (1.6pt);
        \node[caption] at (1.5,1.95) {$\varphi > 0$ inside};
      \end{tikzpicture}\\[4pt]
      {\small\itshape\color{neutral} 1D: zero at the ends}
    \end{column}
  \end{columns}
  \takeaway{The fixed cutoff $\varphi$ pins the boundary exactly and for free;
  the trainable network does all the real work inside.}
\end{frame}

% -----------------------------------------------------------------------------
\begin{frame}{First, read the error as a sum of sine waves}
  We will pick the loss to drive the residual down. Does that drive the
  \emph{error} down too? To find out, write both in the natural basis for
  $-\Delta$ on a box: its \important{sine modes}.
  \smallskip
  \begin{columns}[c]
    \begin{column}{0.57\textwidth}
      \begin{itemize}\itemsep0.5em
        \item Any function vanishing on $\partial\Omega$ is a sum of \important{sine
              modes} $\phi_k = \sqrt{2}\,\sin(k\pi x)$, with $\int \phi_k^2 = 1$:
              $\varepsilon(x) = \sum_k a_k \phi_k(x)$, amplitude $a_k$ per mode.
        \item They are the \important{eigenfunctions} of the Laplacian,
              $-\Delta \phi_k = (k\pi)^2 \phi_k$.
        \item So $-\Delta$ does not mix modes: it scales mode $k$ by $(k\pi)^2$,
              and we can reason one mode at a time.
      \end{itemize}
    \end{column}
    \begin{column}{0.41\textwidth}
      \centering
      \begin{tikzpicture}[>=Latex]
        \draw[->] (-0.15,0) -- (3.4,0) node[right, caption] {$x$};
        \draw[thick, accentdark, domain=0:3, samples=100, smooth]
          plot (\x, {0.62*sin(60*\x)});
        \draw[thick, danger, densely dashed, domain=0:3, samples=140, smooth]
          plot (\x, {0.62*sin(180*\x)});
        \node[caption, accentdark] at (2.75,0.5) {$k=1$};
        \node[caption, danger] at (0.5,0.9) {$k=3$};
      \end{tikzpicture}\\[3pt]
      {\small\itshape\color{neutral} higher $k$, higher frequency\\
       ($\phi_k$ is a mode, $\varphi$ the cutoff)}
    \end{column}
  \end{columns}
  \takeaway{In this basis $-\Delta$ is diagonal, scaling mode $k$ by $(k\pi)^2$.
  That single fact drives both the trap and its fix.}
\end{frame}

% -----------------------------------------------------------------------------
\begin{frame}{The trap: a small residual is not a small error}
  Back to the loss. The obvious choice squares the pointwise residual at sample
  points, $\frac{1}{N}\sum_i |f(x_i) + \Delta u_\theta(x_i)|^2$, which is a
  \important{PINN}. See what it does to our modes.
  \medskip
  Write the error $\varepsilon = u_\theta - u^*$. Since $-\Delta u^* = f$, the
  residual \emph{is} the error through the operator, $r = -\Delta\varepsilon$, so
  it scales mode $k$ by $(k\pi)^2$:
  \[
    \varepsilon = a\,\sin(k\pi x)
    \quad\Longrightarrow\quad
    \norm{\varepsilon}_{L^2} \propto |a|,
    \qquad
    \norm{r}_{L^2} = (k\pi)^2\,|a| .
  \]
  \begin{itemize}\itemsep0.3em
    \item \important{Same error, wildly different residual:} at $k = 1$ versus
          $k = 10$ the error is identical, yet the residual differs by $100\times$.
  \end{itemize}
  \takeaway{The pointwise residual over-weights high-frequency error by $(k\pi)^2$,
  so making it small is not the same as making the error small. We need a norm
  that corrects for this.}
\end{frame}

% =============================================================================
\section{The Deep Fourier Residual method}
% =============================================================================

\begin{frame}{The DFR fix: the dual norm of the weak residual}
  Taylor, Pardo and Muga (2023): stop scoring the pointwise residual; measure it
  in the \important{$H^{-1}$ dual norm} instead (weak form in Part 3).
  \begin{itemize}\itemsep0.5em
    \item We saw these sine modes diagonalize $-\Delta$. The dual norm simply
          \emph{reweights} each mode of the residual by $1/\lambda_k$,
          \[
            \norm{\Res(u_\theta)}_{V^*}^2 = \sum_k \frac{|\hat{\Res}_k|^2}{\lambda_k},
            \qquad \lambda_k = \sum_{i=1}^{d} (\pi k_i / L_i)^2 ,
          \]
          with coefficients $\hat{\Res}_k$ from a \important{Discrete Sine
          Transform} (DST).
    \item $k = (k_1, \dots, k_d)$ is a \important{multi-index}, one frequency per
          axis; $i$ runs over the $d$ axes, of side lengths $L_i$. So in
          $d \ge 2$ the weight couples the axes by a \emph{sum}, not a product.
    \item \good{The point:} on the unit interval of the last two slides
          ($d = 1$, $L = 1$), $\lambda_k = (k\pi)^2$; dividing by it undoes the
          amplification behind the trap, so $\norm{\Res}_{V^*}$ tracks the
          energy error.
  \end{itemize}
\end{frame}

% -----------------------------------------------------------------------------
\begin{frame}{What that reweighting does to training}
  In modes, $\varepsilon = \sum_k a_k \phi_k$ gives $\hat{\Res}_k = \lambda_k a_k$,
  so the two losses score the \emph{same} error very differently:
  \[
    \mathcal{L}_{\text{PINN}} = \sum_k |\hat{\Res}_k|^2 = \sum_k \lambda_k^2 a_k^2,
    \qquad
    \mathcal{L}_{\text{DFR}} = \sum_k \frac{|\hat{\Res}_k|^2}{\lambda_k}
      = \sum_k \lambda_k a_k^2 = \norm{\nabla \varepsilon}_{L^2}^2 .
  \]
  \begin{itemize}\itemsep0.4em
    \item \important{The gradient is damped mode by mode:}
          $\nabla_\theta \mathcal{L}_{\text{DFR}}
           = 2\sum_k \lambda_k^{-1} \hat{\Res}_k \nabla_\theta \hat{\Res}_k$,
          so high frequencies no longer dominate the update.
    \item \important{It preconditions the problem:} in the $a_k$ the Hessian is
          $\mathrm{diag}(\lambda_k)$, not $\mathrm{diag}(\lambda_k^2)$; over $K$
          modes per axis the conditioning goes from $O(K^4)$ to $O(K^2)$.
    \item \alert{The price:} the DST needs a structured grid, so training is
          full-batch on one fixed quadrature grid, on a box, with sines as
          eigenfunctions.
  \end{itemize}
  \takeaway{The DFR loss \emph{is} the energy error, so the objective is better
  shaped, not solved: Part 4 finds plain DFR still not converged at its budget.}
\end{frame}

% -----------------------------------------------------------------------------
\begin{frame}{DFR's guarantee}
  \begin{block}{Error-loss equivalence (Taylor, Pardo, Muga 2023)}
    For a well-posed problem there are constants $\gamma, M > 0$ such that, for
    \emph{any} trial function $w$ (trained or not, so in particular our
    $u_\theta$),
    \[
      \frac{1}{M}\,\norm{\Res(w)}_{V^*}
        \;\le\; \norm{u^* - w}_U
        \;\le\; \frac{1}{\gamma}\,\norm{\Res(w)}_{V^*} .
    \]
  \end{block}
  \medskip
  \begin{itemize}\itemsep0.5em
    \item \good{Two-sided:} the loss bounds the error above \emph{and} below.
    \item Drive the DFR loss down, and the energy-norm error goes down with it.
    \item We prove this, and everything that follows, in Part 3.
  \end{itemize}
\end{frame}

% =============================================================================
\section{Quantities of interest}
% =============================================================================

\begin{frame}{Global error is often not what we want}
  DFR gives tight, two-sided control of the \emph{whole-field} energy error. Often
  that is more than we need: the object we actually care about is one
  \important{functional} of the solution, a quantity of interest (QoI):
  \medskip
  \begin{center}
  \begin{tikzpicture}[node distance=4mm and 12mm]
    \node[opnode] (pt) {a point value\\$u^*(x_0)$};
    \node[opnode, right=of pt] (avg) {a subdomain\\average};
    \node[opnode, right=of avg] (flux) {a boundary\\flux};
  \end{tikzpicture}
  \end{center}
  \medskip
  \begin{itemize}\itemsep0.5em
    \item A small global error is \emph{sufficient} but not \emph{necessary} for
          an accurate QoI.
    \item Driving the global error to a tight tolerance is wasteful when only
          $J(u^*)$ matters.
  \end{itemize}
  \takeaway{We want a loss aimed at the quantity of interest, not the whole
  field. That is Part 2.}
\end{frame}

\end{document}
